\chapter{Висновки}

У результаті проведеного дослідження:

\begin{enumerate}
    \item Зібрано та систематизовано дані про 132 акредитовані освітні програми
    зі спеціальності 022 «Дизайн» в Україні.

    \item Проведено статистичний аналіз структури навчальних планів:
    \begin{itemize}
        \item Середня кількість освітніх компонентів на програму складає близько 30
        \item Співвідношення обов'язкової та вибіркової частин відносно стабільне
        \item Бакалаврські програми (74 од.) домінують у загальній структурі
    \end{itemize}

    \item За допомогою тематичного моделювання виявлено основні тематичні кластери:
    графічний дизайн, UX/UI, 3D-моделювання, традиційні мистецтва.

    \item Аналіз прогалин показав недостатнє покриття AI-навичок:
    \begin{itemize}
        \item Покриття базових AI-навичок: менше 20\%
        \item Покриття цифрових інструментів: близько 60\%
        \item Генеративний AI практично відсутній у навчальних планах
    \end{itemize}

    \item Розроблено пропозиції щодо модернізації навчальних програм
    з інтеграцією AI-інструментів.
\end{enumerate}

\section{Перспективи подальших досліджень}

\begin{itemize}
    \item Порівняльний аналіз із зарубіжними програмами
    \item Моніторинг впровадження AI у дизайн-освіту
    \item Опитування випускників та роботодавців
\end{itemize}
